
Dado el siguiente problema de programación lineal:

\begin{equation}
\begin{split}
\mbox{max. } z = 3x_1 + 2x_2\\
s.a.:\\
2x_1 + 4x_2 \leq 80\\
4x_1 + 3x_2 \leq 60\\
 x_1 +  x_2 \leq 15\\
x_1,\,\,x_2\geq 0\\
\end{split}
\end{equation}

Se conoce la solución óptima.

\begin{center}
\begin{tabular}{c|c|ccccc|}
 & $z$ & $x_1$ & $x_2$ & $h_1$ & $h_2$ & $h_3$\\ 
 & -45 & 0 & -1/4 & 0 & -3/4 & 0 \\ 
\hline
$h_1$ & 50  & 0 & 5/2 & 1 & -1/2 & 0\\ 
$x_1$ & 15  & 1 & 3/4 & 0 & 1/4 & 0\\ 
$h_3$ & 0  & 0 & 1/4 & 0 & -1/4 & 1\\ 
\hline
\end{tabular}
\end{center}



\begin{enumerate}
    \item Se pide formular el problema dual asociado.

\item Dada la siguiente tabla, que corresponde a una solución óptima del problema dual asociado, comprobar que se cumple el teorema fundamental de la dualidad entre el problema primal y el problema dual y, sus respectivas soluciones óptimas. 

\begin{center}
\begin{tabular}{c|c|ccccc|}
 & $z$ & $y_1$ & $y_2$ & $y_3$ & $h_1$ & $h_2$\\ 
 & 45 & -50 & 0 & 0 & -15 & 0 \\ 
\hline
$h_2$ & 1  & -2 & 1 & 0 & -1 & 1\\ 
$y_3$ & 3  & 2 & 4 & 1 & -1 & 0\\ 
\hline
\end{tabular}
\end{center}

\end{enumerate}     